\section{Finite gauge, horizon, and dense matter projections}

\begin{proposition}[Finite gauge field carrier]\label{thm:finite-gauge-group}\lean{Obsmetrologia.DenseMatter.Quark.FiniteGroup}\leanok\uses{thm:k4-group}
A finite gauge group is enough to state discrete gauge transformations, path
holonomy, and Wilson-loop invariance without pretending to have a Lie group or
continuum bundle.
\end{proposition}

\begin{theorem}[Wilson loops are gauge-invariant class functions]\label{thm:wilson-loop-invariant}\lean{Obsmetrologia.DenseMatter.Quark.wilsonLoop_invariant}\leanok\uses{thm:finite-gauge-group}
A class function of a closed-loop holonomy is invariant under finite gauge
transformation.  This is the finite skeleton of the Wilson-loop observable; the
analytic area law, spinors, and Lie-bundle layer remain outside the present
formalization.
\end{theorem}

\begin{theorem}[The $\mathbb{Z}/2$ flux reading: contextuality and the cloth mirror are gauge statements]\label{thm:z2-flux-reading}\lean{Obsmetrologia.DenseMatter.Quark.z2_flux_reading}\leanok\uses{thm:wilson-loop-invariant, thm:pr-cohomological-contextuality}
Three welds through the $\mathbb{Z}/2$ gauge group. (i) \emph{Pure gauge is the
coboundary condition}: a two-argument Bool target is the
\texttt{QuantumContextuality.delta0} gradient of a global assignment iff its
connection on the complete edge set is the gauge transform of the trivial
connection (\texttt{pureGauge\_iff\_delta0}); a pure-gauge connection has trivial
Wilson loops (\texttt{pureGauge\_closed\_holo\_one}). (ii) \emph{The PR box
carries one unit of flux}: its parity target, read as a $\mathbb{Z}/2$ connection
on the four measurements, has Wilson loop \texttt{true} around the context
4-cycle and is not pure gauge (\texttt{prConnection\_flux}) --- strong
contextuality as gauge curvature, the gauge-theoretic re-proof of
\cref{thm:pr-cohomological-contextuality}. (iii) \emph{The cloth mirror is a
parity law}: flipping every edge bit --- the verso reading, the over/under swap
that acts by bilattice negation (\texttt{overUnder\_swap\_is\_negate}) ---
changes a loop's holonomy by the parity of its length
(\texttt{holo\_neg}): even loops are face-independent
(\texttt{verso\_even\_loop\_invariant}), odd loops flip
(\texttt{verso\_odd\_loop\_flips}) --- the M\"obius seam of
\cref{thm:mobius-band-holonomy} as the odd case. The bundle depends on no
axioms: the PR measurement carrier is the four-corner cell
`Bool × Bool` (party × setting, \texttt{PRM}), so the flux route is the
axiom-free proof of PR contextuality.
\end{theorem}

\begin{theorem}[The product connection: one loop, two charges]\label{thm:product-connection}\lean{Obsmetrologia.DenseMatter.Quark.product_connection_reading}\leanok\uses{thm:z2-flux-reading, thm:wilson-loop-invariant}
The product of two finite gauge groups acts componentwise
(\texttt{FiniteGroup.prod}), and the holonomy of a product connection is the
pair of the factor holonomies (\texttt{holo\_prod}): one Wilson loop reads
both charges at once. Instantiated on the PR measurement carrier with the
contextuality flux in one factor and the orientation (length-parity)
connection in the other (\texttt{orientation\_holo}), the context 4-cycle's
product Wilson loop reads \((\mathrm{true}, \mathrm{false})\) --- one unit
of contextual flux, even orientation --- while any odd loop flips the second
factor. The whole development depends on no axioms. The continuation named
here is discharged by \cref{thm:curvature-contextuality-one-loop}.
\end{theorem}

\begin{theorem}[Curvature and contextuality are one gauge charge]\label{thm:curvature-contextuality-one-loop}\lean{Obsmetrologia.DenseMatter.Quark.curvature_contextuality_one_loop}\leanok\uses{thm:product-connection, thm:disclination-holonomy-quantum, thm:z2-flux-reading}
The hexagonal frame group as a gauge group: \texttt{C6} is the rot-orbit
of $(1,0)$ in the Eisenstein plane under Eisenstein multiplication --- the
sixth roots of unity as a derived carrier, no posited finite type --- and
on the orbit, multiplication by the generator is the disclination
carrier's derived rotation (\texttt{zeta\_mul\_is\_rot\_on\_orbit}). The
constant one-sixth-turn connection reads path length as rotational
holonomy (\texttt{zetaStep\_holo}); four- and five-loops carry a nonzero
frame deficit while the six-loop closes flat
(\texttt{c6\_deficit\_spectrum}). The unified connection pairs this
spatial factor with the PR connection on the PR measurement sites: its
single Wilson loop around the context 4-cycle reads
$(\zeta^4, \mathrm{true})$ (\texttt{unified\_wilson}) --- a nonzero
curvature charge and the contextual flux, one loop, one product charge,
two obstructions --- gauge-invariantly, by
\cref{thm:wilson-loop-invariant}. The whole development depends on no
axioms.
\end{theorem}

\begin{theorem}[The loop-length dictionary and the deficit structure]\label{thm:loop-length-dictionary}\lean{Obsmetrologia.DenseMatter.Quark.deficit_charge_structure}\leanok\uses{thm:curvature-contextuality-one-loop}
The generator of the hexagonal frame group has order six as a power law
(\texttt{c6pow\_eq\_one\_iff}: a power of the unit turn is the identity
exactly at multiples of six), so on the constant one-sixth-turn fabric a
Wilson loop is flat exactly when its length is divisible by six
(\texttt{zetaStep\_wilson\_iff}) --- every context cycle of every
measurement scenario is a curvature census entry, its deficit read off
its length. The deficit law (\texttt{c6pow\_inv\_deficit}): an
$L$-loop's charge is the inverse of its $(6-L)$-quantum deficit. The
structure of the PR charge (\texttt{pr\_deficit\_structure}): the
unified loop's curvature factor $\zeta^4$ is the inverse \emph{square}
of the unit turn --- the PR plaquette is short two coordinations --- and
the square of the generator satisfies the Eisenstein relation
$\zeta^2 = \zeta - 1$. The whole development depends on no axioms.
\end{theorem}

\begin{theorem}[Holonomy factors over path concatenation]\label{thm:holo-append}\lean{Obsmetrologia.DenseMatter.Quark.holo_append}\leanok\uses{thm:wilson-loop-invariant}
The front/back splitting law any cup-product-shaped construction on paths must go through: for a
single gauge group $G$, the holonomy along a concatenated path $u \mathbin{+\kern-0.6em+} v$ is the
product of the front sub-path's holonomy and the back sub-path's, the back one read from the
front's endpoint (\texttt{holo\_append}, by induction peeling one edge and reassociating; depends
on no axioms).
\end{theorem}

\begin{theorem}[The unified Wilson loop is not a front/back cup product]\label{thm:concrete-cup-product}\notready\uses{thm:curvature-contextuality-one-loop, thm:z2-flux-reading, thm:holo-append}
A candidate route to a cup-product reading of \cref{thm:curvature-contextuality-one-loop}, tested
and found not to work as stated. The unified connection \texttt{unifiedPRConnection} reads both
the $\mathbb{Z}/2$ contextual class and the $\mathbb{Z}/6$ deficit class on \emph{every} edge of
the PR 4-cycle simultaneously (\texttt{holo\_prod} splits by \emph{coefficient factor}, the same
full loop read through two group projections); an Alexander--Whitney cup product would instead
split by \emph{path position} (\cref{thm:holo-append}'s front/back), reading one carrier on the
front edges and the other on the back. These are different operations, and a direct computation
confirms they disagree here: embedding $\zeta$ on a front sub-path and \texttt{prConnection} on
the complementary back sub-path, for every split point of the 4-cycle, never reproduces the target
charge $(\zeta^4,\mathrm{true})$ --- the front-only holonomy through $\zeta$ needs all four edges
to reach $\zeta^4$, leaving no edges for the $\mathbb{Z}/2$ back half to read a nontrivial
contextual class from. \texttt{unifiedPRConnection} is not itself split into a
$\mathbb{Z}/2$-only front and $\mathbb{Z}/6$-only back at any point, so no re-derivation of
\cref{thm:curvature-contextuality-one-loop} as a genuine front/back cup evaluation exists on this
connection; a cup-product reading, if one is provable at all, needs a different carrier where the
two coefficient systems occupy disjoint sub-paths rather than every edge at once.
\end{theorem}

\begin{theorem}[Exchange as graded anticommutativity]\label{thm:exchange-graded-commutativity}\notready\uses{thm:concrete-cup-product}
Once a genuine cup-product carrier for \cref{thm:concrete-cup-product} exists: cup products of
degree-one classes are graded-anticommutative in the general theory
($\alpha\cup\beta = -(\beta\cup\alpha)$), so trading the order of the contextual and deficit
charges would cost a sign. Stating and proving this needs a coefficient ring admitting the sign
in the first place --- the $\mathbb{Z}/2$ and $\mathbb{Z}/6$ factors do not obviously carry one
jointly --- so this remains downstream of finding that carrier, not an independent claim.
\end{theorem}

\begin{theorem}[Finite SO readout-frame law]\label{thm:finite-so-readout-frame}\lean{Obsmetrologia.CausalSets.BlackHoles.Horizon.FiniteSOReadoutFrame, Obsmetrologia.CausalSets.BlackHoles.Horizon.FiniteSOReadoutFrame.horizon_surface_eq_shroud, Obsmetrologia.CausalSets.BlackHoles.Horizon.FiniteSOReadoutFrame.horizon_interior_eq_blackHoleRegion, Obsmetrologia.CausalSets.BlackHoles.Horizon.FiniteSOReadoutFrame.horizon_exterior_eq_readoutExterior, Obsmetrologia.CausalSets.BlackHoles.Horizon.FiniteSOReadoutFrame.closedLoop_holonomy_conjugates, Obsmetrologia.CausalSets.BlackHoles.Horizon.FiniteSOReadoutFrame.wilsonLoop_invariant}\leanok\uses{thm:finite-gauge-group, thm:wilson-loop-invariant, thm:readout-relative-horizon-frontier-laws}
The SO item is carried only as a finite readout-frame law.  A
\texttt{FiniteSOReadoutFrame} pairs a readout-generated horizon with a finite
gauge group used as the orientation-frame proxy, a finite connection, a local
gauge transform, and a closed loop.  The horizon side still computes its
surface, interior, and exterior from the readout; the frame side inherits the
closed-loop holonomy conjugation and Wilson class-function invariance already
proved for finite gauge groups.  No smooth \(\mathrm{SO}(n)\), determinant
condition, frame bundle, spin structure, or continuum Rindler metric is
asserted.
\end{theorem}

\begin{proposition}[Finite hole no-readout region]\label{thm:hole-region}\lean{Obsmetrologia.CausalSets.BlackHoles.HoleRegion}\leanok\uses{thm:causal-settling-transitive-closure, thm:cone-field-geometry}
Given a causal set and a chosen readout region, the hole region is the
no-readout region: events that do not causally precede any readout event.
\end{proposition}

\begin{theorem}[No-readout is future-closed]\label{thm:hole-future-closed}\lean{Obsmetrologia.CausalSets.BlackHoles.hole_future_closed}\leanok\uses{thm:hole-region}
The finite hole region is future-closed: once an event cannot reach readout,
no causal future of it can reach readout either.
\end{theorem}

\begin{theorem}[The finite no-readout region is a causal void]\label{thm:causal-void}\lean{Obsmetrologia.CausalSets.BlackHoles.CausalVoid.no_future_to_exterior}\leanok\uses{thm:hole-region, thm:hole-future-closed}
A \texttt{CausalVoid} packages the generated no-readout interior with its readout
exterior.  Its void law is not absence by metaphor: no interior event causally
precedes a readout-exterior event.
\end{theorem}

\begin{proposition}[Finite hole shroud]\label{thm:hole-shroud}\lean{Obsmetrologia.CausalSets.BlackHoles.Shroud}\leanok\uses{thm:causal-void}
The shroud is the no-readout side of the boundary: a hole-region event
causally adjacent to a readout-exterior event.
\end{proposition}

\begin{proposition}[Shroud-indexed keratin]\label{thm:shroud-indexed-keratin}\lean{Obsmetrologia.AperiodicOrder.MaterialSectors.Contact.ShroudIndexedKeratin}\leanok\uses{thm:hole-shroud}
A shroud-indexed keratin carrier welds a hole shroud to directed contact
and irreversible lock history without asserting a disc, entropy law, or smooth
horizon theorem.
\end{proposition}

\begin{theorem}[Finite black-hole-region proxy]\label{thm:black-hole-region-proxy}\lean{Obsmetrologia.CausalSets.BlackHoles.Horizon.BlackHoleRegion, Obsmetrologia.CausalSets.BlackHoles.Horizon.blackHole_future_closed}\leanok\uses{thm:hole-region, thm:hole-future-closed}
The external black-hole-region formula is represented here only in the finite
readout-relative sense: \texttt{BlackHoleRegion} is the already proved no-readout
\texttt{HoleRegion}, hence it is future-closed.  The name does not add future
null infinity, smooth null generators, expansion scalars, or continuum
singularity theorems.
\end{theorem}

\begin{theorem}[Finite apparent event and plate laws]\label{thm:finite-apparent-event-plate}\lean{Obsmetrologia.CausalSets.BlackHoles.Horizon.ApparentEvent, Obsmetrologia.CausalSets.BlackHoles.Horizon.ApparentEvent.in_blackHoleRegion, Obsmetrologia.CausalSets.BlackHoles.Horizon.ApparentEvent.not_exterior, Obsmetrologia.CausalSets.BlackHoles.Horizon.ApparentEvent.has_exterior_predecessor, Obsmetrologia.CausalSets.BlackHoles.Horizon.BekensteinPlate, Obsmetrologia.CausalSets.BlackHoles.Horizon.mem_bekensteinPlate_iff, Obsmetrologia.CausalSets.BlackHoles.Horizon.bekensteinPlate_in_blackHoleRegion}\leanok\uses{thm:hole-shroud, thm:black-hole-region-proxy}
An \texttt{ApparentEvent} is a finite shroud event: it lies in the no-readout
side, is not exterior, and has an exterior predecessor.  The
\texttt{BekensteinPlate} is only the finite list of shroud events; cardinality or
area/capacity laws remain to be proved over that list.
\end{theorem}

\begin{theorem}[Finite Bekenstein plate counting law]\label{thm:finite-bekenstein-plate-counting}\lean{Obsmetrologia.CausalSets.BlackHoles.Horizon.bekensteinPlate_card_eq_filter, Obsmetrologia.CausalSets.BlackHoles.Horizon.bekensteinPlate_length_le_events}\leanok\uses{thm:finite-apparent-event-plate}
The first counting statement proved over the finite Bekenstein plate is just
the list law already carried by the shroud-event carrier: the plate cardinality
is the length of the shroud-filtered event list, and is bounded above by the
ambient finite causal-set event list.  This is a finite counting bound over
events, not an area law, entropy law, or continuum Bekenstein bound.
\end{theorem}

\begin{theorem}[Finite Bekenstein plate capacity witness]\label{thm:finite-bekenstein-plate-capacity}\lean{Obsmetrologia.CausalSets.BlackHoles.Horizon.BekensteinPlateCapacityWitness, Obsmetrologia.CausalSets.BlackHoles.Horizon.BekensteinPlateCapacityWitness.plate, Obsmetrologia.CausalSets.BlackHoles.Horizon.BekensteinPlateCapacityWitness.capacity, Obsmetrologia.CausalSets.BlackHoles.Horizon.BekensteinPlateCapacityWitness.capacity_pos, Obsmetrologia.CausalSets.BlackHoles.Horizon.BekensteinPlateCapacityWitness.capacity_eq_landauer_cost, Obsmetrologia.CausalSets.BlackHoles.Horizon.BekensteinPlateCapacityWitness.ledger_eq_eulerChar, Obsmetrologia.CausalSets.BlackHoles.Horizon.BekensteinPlateCapacityWitness.plate_card_eq_filter, Obsmetrologia.CausalSets.BlackHoles.Horizon.BekensteinPlateCapacityWitness.plate_length_le_events, Obsmetrologia.CausalSets.BlackHoles.Horizon.BekensteinPlateCapacityWitness.event_in_blackHoleRegion, Obsmetrologia.CausalSets.BlackHoles.Horizon.BekensteinPlateCapacityWitness.finite_capacity_package}\leanok\uses{thm:finite-bekenstein-plate-counting, thm:load-bearing-landauer-support}
A finite Bekenstein plate capacity witness pairs the shroud-event plate with a
load-bearing Landauer support.  Its capacity is the positive finite Landauer
cost of that support, while the plate side still has exactly the shroud-filter
count and ambient event-list bound.  The witness also exposes the existing
anti-Landauer ledger equality to Euler characteristic.  This is a finite
information-cost package over the plate carrier, not an entropy formula,
continuum area law, or Bekenstein bound.
\end{theorem}

\begin{theorem}[Finite horizon membrane dynamics]\label{thm:finite-horizon-membrane-dynamics}\lean{Obsmetrologia.CausalSets.BlackHoles.Horizon.HorizonMembrane.vertex_in_interior, Obsmetrologia.CausalSets.BlackHoles.Horizon.HorizonMembrane.lockedSeed_infl, Obsmetrologia.CausalSets.BlackHoles.Horizon.HorizonMembrane.settle_locked_least_compatible}\leanok\uses{thm:shroud-indexed-keratin, thm:keratin-settle-lock}
The finite membrane object is the shroud-indexed keratin carrier read at the
horizon.  Besides the shroud-address law, its addressed vertices lie in the
hole interior, its lock histories inflate every seed stalk, and settling the
locked seed is least among compatible states above it.  This is a finite
contact-and-settlement law for a shroud-indexed membrane; it is not a continuum
membrane paradigm with surface stress, viscosity, conductivity, or transport
coefficients.
\end{theorem}

\begin{theorem}[Readout-relative horizon, membrane, patch, and censorship laws]\label{thm:readout-relative-horizon-frontier-laws}\lean{Obsmetrologia.CausalSets.BlackHoles.Horizon.ReadoutRelativeHorizon, Obsmetrologia.CausalSets.BlackHoles.Horizon.readoutRelativeHorizon_surface_interior, Obsmetrologia.CausalSets.BlackHoles.Horizon.HorizonMembrane, Obsmetrologia.CausalSets.BlackHoles.Horizon.HorizonMembrane.vertex_on_shroud, Obsmetrologia.CausalSets.BlackHoles.Horizon.HoleRegionPatch, Obsmetrologia.CausalSets.BlackHoles.Horizon.shroudPatch, Obsmetrologia.CausalSets.BlackHoles.Horizon.no_naked_hole_effect_source}\leanok\uses{thm:hole-shroud, thm:shroud-indexed-keratin, thm:finite-apparent-event-plate}
The Rindler-style lesson is carried only as readout relativity: a chosen readout
generates a horizon package.  The membrane carrier is the existing
shroud-indexed keratin object; its vertices index shroud points.  The shroud is
a hole-region patch, and the finite censorship law says a hole effect cannot
have an exterior source.
\end{theorem}

\begin{theorem}[Finite readout patch and censorship consolidation]\label{thm:finite-readout-patch-censorship-consolidation}\lean{Obsmetrologia.CausalSets.BlackHoles.Horizon.readoutRelativeHorizon_surface_eq_shroud, Obsmetrologia.CausalSets.BlackHoles.Horizon.readoutRelativeHorizon_interior_eq_blackHoleRegion, Obsmetrologia.CausalSets.BlackHoles.Horizon.readoutRelativeHorizon_exterior_eq_readoutExterior, Obsmetrologia.CausalSets.BlackHoles.Horizon.readoutRelativeHorizon_surface_not_exterior, Obsmetrologia.CausalSets.BlackHoles.Horizon.readoutRelativeHorizon_surface_has_exterior_predecessor, Obsmetrologia.CausalSets.BlackHoles.Horizon.antihole_no_listed_readout_ingress, Obsmetrologia.CausalSets.BlackHoles.Horizon.listed_readout_not_antihole, Obsmetrologia.CausalSets.BlackHoles.Horizon.HoleRegionPatch.in_blackHoleRegion, Obsmetrologia.CausalSets.BlackHoles.Horizon.HoleRegionPatch.not_exterior, Obsmetrologia.CausalSets.BlackHoles.Horizon.HoleRegionPatch.no_future_to_exterior, Obsmetrologia.CausalSets.BlackHoles.Horizon.shroudPatch_patch_eq_shroud, Obsmetrologia.CausalSets.BlackHoles.Horizon.shroudPatch_has_exterior_predecessor, Obsmetrologia.CausalSets.BlackHoles.Horizon.shroudPatch_not_exterior, Obsmetrologia.CausalSets.BlackHoles.Horizon.kinematicHoleRegionPatch, Obsmetrologia.CausalSets.BlackHoles.Horizon.kinematicHoleRegionPatch_patch_eq_holeRegion, Obsmetrologia.CausalSets.BlackHoles.Horizon.kinematicHoleRegionPatch_not_exterior, Obsmetrologia.CausalSets.BlackHoles.Horizon.kinematicHoleRegionPatch_no_future_to_exterior, Obsmetrologia.CausalSets.BlackHoles.Horizon.no_naked_canonical_hole_effect_source}\leanok\uses{thm:readout-relative-horizon-frontier-laws, thm:antihole-readout-complementarity}
The readout-relative horizon package is now tied back to its computed carrier:
its surface is the shroud, its interior is the no-readout black-hole-region
proxy, its exterior is the readout exterior, and surface points have exterior
predecessors while not being exterior themselves.  Antiholes admit no ingress
from listed readout anchors, and listed readout points are not antihole points.
For arbitrary hole-region patches, membership gives black-hole-region
membership, rules out readout-exterior membership, and prevents causal futures
to exterior points.  The shroud patch is definitionally the shroud and inherits
the exterior-predecessor law.  The same patch law now has a named kinematic
specialization: on the sequencing causal-diamond space \(K(C)\), the no-readout
predicate itself is a hole-region patch and has the same not-exterior and
no-future-to-exterior laws.  For the canonical hole generated by a readout,
the finite censorship law is the concrete no-naked-source statement: an effect
source is not in the readout exterior.
\end{theorem}

\begin{theorem}[Finite antihole and readout-complementarity laws]\label{thm:antihole-readout-complementarity}\lean{Obsmetrologia.CausalSets.BlackHoles.Horizon.ReadoutFuture, Obsmetrologia.CausalSets.BlackHoles.Horizon.AntiholeRegion, Obsmetrologia.CausalSets.BlackHoles.Horizon.readoutFuture_future_closed, Obsmetrologia.CausalSets.BlackHoles.Horizon.antihole_past_closed, Obsmetrologia.CausalSets.BlackHoles.Horizon.antihole_no_ingress_from_readoutFuture, Obsmetrologia.CausalSets.BlackHoles.Horizon.AntiholeShroud, Obsmetrologia.CausalSets.BlackHoles.Horizon.antiholeShroud_has_readoutFuture_successor, Obsmetrologia.CausalSets.BlackHoles.Horizon.ReadoutComplementarityWitness, Obsmetrologia.CausalSets.BlackHoles.Horizon.ReadoutComplementarityWitness.not_single_readout}\leanok\uses{thm:hole-region, thm:black-hole-region-proxy}
The finite antihole is the order-dual no-ingress side of a readout: events not
reached from the chosen readout.  The readout future is future-closed, the
antihole region is past-closed, no readout-future event can causally precede an
antihole event, and an antihole shroud touches the readout future only by a
successor.  Readout complementarity is proved only as a finite witness: a point
that is a hole for one readout and exterior for another proves those two
readouts cannot be collapsed to one readout.
\end{theorem}

\begin{theorem}[Finite readout no-comparison law]\label{thm:finite-readout-no-comparison}\lean{Obsmetrologia.CausalSets.BlackHoles.Horizon.ReadoutComplementarityWitness.left_not_exterior, Obsmetrologia.CausalSets.BlackHoles.Horizon.ReadoutComplementarityWitness.right_not_blackHoleRegion, Obsmetrologia.CausalSets.BlackHoles.Horizon.ReadoutComplementarityWitness.no_common_classification}\leanok\uses{thm:antihole-readout-complementarity}
A finite readout-complementarity witness is not merely evidence that two
readout predicates differ.  At the witnessed event, the left readout's hole
classification rules out left exterior membership, the right readout's exterior
classification rules out right hole membership, and the two readouts cannot
share the same hole/exterior classification pair at that event.  This is a
pointwise no-comparison law over finite readout classifications, not a continuum
black-hole-complementarity theorem.
\end{theorem}

\begin{theorem}[Finite wormhole bridge witness]\label{thm:finite-wormhole-bridge}\lean{Obsmetrologia.CausalSets.BlackHoles.Horizon.FiniteWormholeBridge, Obsmetrologia.CausalSets.BlackHoles.Horizon.FiniteWormholeBridge.causal_bridge, Obsmetrologia.CausalSets.BlackHoles.Horizon.FiniteWormholeBridge.not_factorizes, Obsmetrologia.CausalSets.BlackHoles.Horizon.eprFiniteWormholeBridge}\leanok\uses{thm:molecular-chaos-factorization, thm:readout-relative-horizon-frontier-laws}
The finite wormhole bridge is not a continuum throat.  It is a causal adjacency
equipped with a nonzero connected-correlation minor from the kinetic gluing
module.  The bridge therefore has a causal edge in the finite spacetime carrier,
and its joint data cannot be molecular-chaos factorized.  The concrete EPR
joint supplies a non-vacuous bridge whenever the causal carrier supplies an
adjacent event pair.
\end{theorem}

\begin{theorem}[Finite geon carrier]\label{thm:finite-geon-carrier}\lean{Obsmetrologia.CausalSets.BlackHoles.Horizon.FiniteGeonCarrier, Obsmetrologia.CausalSets.BlackHoles.Horizon.FiniteGeonCarrier.futureAnchor, Obsmetrologia.CausalSets.BlackHoles.Horizon.FiniteGeonCarrier.futureCore, Obsmetrologia.CausalSets.BlackHoles.Horizon.FiniteGeonCarrier.connected_region, Obsmetrologia.CausalSets.BlackHoles.Horizon.FiniteGeonCarrier.topologicalRegion, Obsmetrologia.CausalSets.BlackHoles.Horizon.FiniteGeonCarrier.shroud_between_exterior_and_root}\leanok\uses{thm:hole-region, thm:finite-wormhole-bridge}
A finite geon carrier is a finite hole-side settling process whose stages reach
one common no-readout root.  The carrier is therefore self-sustaining only in
the local sense already proved by the hole-region development: it induces a
future anchor, a connected future core, a connected computed hole region, and a
certified Alexandrov topological region.  Shroud events lie between exterior
predecessors and the settled root.  No continuum Wheeler-geon geometry or mass
law is asserted.
\end{theorem}

\begin{theorem}[Finite ergosphere extraction witness]\label{thm:finite-ergosphere-extraction}\lean{Obsmetrologia.CausalSets.BlackHoles.Horizon.FiniteErgosphereExtraction, Obsmetrologia.CausalSets.BlackHoles.Horizon.FiniteErgosphereExtraction.boundary_in_blackHoleRegion, Obsmetrologia.CausalSets.BlackHoles.Horizon.FiniteErgosphereExtraction.has_exterior_predecessor, Obsmetrologia.CausalSets.BlackHoles.Horizon.FiniteErgosphereExtraction.strict_energy_drop}\leanok\uses{thm:finite-apparent-event-plate, thm:free-energy-lyapunov}
The finite ergosphere extraction witness is a boundary-plus-energy witness, not
a rotating continuum region.  It packages a finite apparent event on the shroud
with a hard-closure settling firing on a finite support cell.  The boundary part
places the witness inside the computed black-hole-region proxy and supplies an
exterior predecessor; the energy part proves that the Nat-valued free
energy strictly drops across the firing.  This is the finite sense in which
``work'' is extractable at the horizon frontier.
\end{theorem}

\begin{theorem}[Finite Unruh readout witness]\label{thm:finite-unruh-readout-witness}\lean{Obsmetrologia.CausalSets.BlackHoles.Horizon.FiniteUnruhReadoutWitness, Obsmetrologia.CausalSets.BlackHoles.Horizon.FiniteUnruhReadoutWitness.boundary_in_blackHoleRegion, Obsmetrologia.CausalSets.BlackHoles.Horizon.FiniteUnruhReadoutWitness.has_exterior_predecessor, Obsmetrologia.CausalSets.BlackHoles.Horizon.FiniteUnruhReadoutWitness.sector_particleNumber, Obsmetrologia.CausalSets.BlackHoles.Horizon.FiniteUnruhReadoutWitness.singleCreation, Obsmetrologia.CausalSets.BlackHoles.Horizon.FiniteUnruhReadoutWitness.singleCreation_sector_particleNumber}\leanok\uses{thm:finite-apparent-event-plate, thm:c4-ledger-fock-sector}
The finite Unruh readout witness is only a readout-relative particle-number
witness.  It packages a finite apparent boundary event with a C4
creation/annihilation ledger and its generated Fock sector.  The boundary part
places the event inside the computed black-hole-region proxy and gives an
exterior predecessor; the Fock part proves that the read sector has exactly the
ledger particle number.  The single-creation constructor gives a one-particle
boundary readout.  No thermal bath, detector temperature, or continuum quantum
field theorem is asserted.
\end{theorem}

\begin{theorem}[Finite Unruh thermal-readout bridge]\label{thm:finite-unruh-thermal-readout-bridge}\lean{Obsmetrologia.CausalSets.BlackHoles.Horizon.FiniteUnruhThermalReadoutWitness, Obsmetrologia.CausalSets.BlackHoles.Horizon.FiniteUnruhThermalReadoutWitness.boundary_in_blackHoleRegion, Obsmetrologia.CausalSets.BlackHoles.Horizon.FiniteUnruhThermalReadoutWitness.has_exterior_predecessor, Obsmetrologia.CausalSets.BlackHoles.Horizon.FiniteUnruhThermalReadoutWitness.sector_particleNumber, Obsmetrologia.CausalSets.BlackHoles.Horizon.FiniteUnruhThermalReadoutWitness.creationCount_pos, Obsmetrologia.CausalSets.BlackHoles.Horizon.FiniteUnruhThermalReadoutWitness.flux_balanced, Obsmetrologia.CausalSets.BlackHoles.Horizon.FiniteUnruhThermalReadoutWitness.temperature_eq_distribution, Obsmetrologia.CausalSets.BlackHoles.Horizon.FiniteUnruhThermalReadoutWitness.produces_creation_landauer_and_discharge, Obsmetrologia.CausalSets.BlackHoles.Horizon.FiniteUnruhThermalReadoutWitness.readout_and_temperature}\leanok\uses{thm:finite-unruh-readout-witness, thm:finite-thermodynamic-transition-production}
The next finite Unruh carrier pairs the readout-relative particle-number
witness with the already proved thermodynamic transition carrier.  The boundary
and Fock laws remain those of \texttt{FiniteUnruhReadoutWitness}; the
thermodynamic side contributes a positive creation count, balanced stochastic
flux, stationary-distribution temperature, positive Landauer cost, and positive
hard-flow discharge.  This is the finite thermal-readout bridge only.  It still
does not assert a continuum thermal bath, detector-response spectrum, or quantum
field theorem.
\end{theorem}

\begin{theorem}[Finite trapped-surface witness]\label{thm:finite-trapped-surface-witness}\lean{Obsmetrologia.CausalSets.BlackHoles.Horizon.FiniteTrappedSurfaceWitness, Obsmetrologia.CausalSets.BlackHoles.Horizon.FiniteTrappedSurfaceWitness.crosses_readout_horizon, Obsmetrologia.CausalSets.BlackHoles.Horizon.FiniteTrappedSurfaceWitness.lower_exterior, Obsmetrologia.CausalSets.BlackHoles.Horizon.FiniteTrappedSurfaceWitness.upper_shroud, Obsmetrologia.CausalSets.BlackHoles.Horizon.FiniteTrappedSurfaceWitness.upper_in_blackHoleRegion, Obsmetrologia.CausalSets.BlackHoles.Horizon.FiniteTrappedSurfaceWitness.upper_not_exterior, Obsmetrologia.CausalSets.BlackHoles.Horizon.FiniteTrappedSurfaceWitness.apparentEvent, Obsmetrologia.CausalSets.BlackHoles.Horizon.FiniteTrappedSurfaceWitness.upper_has_exterior_predecessor}\leanok\uses{thm:kugelblitz-local-hole-entry, thm:finite-apparent-event-plate}
A finite trapped-surface witness is the kugelblitz articulation site read as a
boundary carrier: the next articulation step supplies a diamond whose lower
endpoint is in the readout exterior and whose upper endpoint lands on the
shroud.  The upper endpoint is therefore inside the computed black-hole-region
proxy, is not exterior, and forms a finite apparent event with an exterior
predecessor.  This proves the exterior-to-shroud trapping direction only;
the continuum trapped-surface singularity theorem and the strict-growth
exactness statement remain the separate kugelblitz growth conjecture.
\end{theorem}

\begin{theorem}[Finite critical-collapse carrier laws]\label{thm:finite-critical-collapse-carrier-laws}\lean{Obsmetrologia.CausalSets.BlackHoles.Horizon.CriticalCollapseCandidate, Obsmetrologia.CausalSets.BlackHoles.Horizon.criticalCollapseCandidate_site, Obsmetrologia.CausalSets.BlackHoles.Horizon.criticalCollapseCandidate_tail_mem_prefixes, Obsmetrologia.CausalSets.BlackHoles.Horizon.criticalCollapseCandidate_next_mem_prefixes, Obsmetrologia.CausalSets.BlackHoles.Horizon.criticalCollapseCandidate_tail_reaches_next, Obsmetrologia.CausalSets.BlackHoles.Horizon.criticalCollapseCandidate_tail_causalLe_next, Obsmetrologia.CausalSets.BlackHoles.Horizon.criticalCollapseCandidate_tail_exterior, Obsmetrologia.CausalSets.BlackHoles.Horizon.criticalCollapseCandidate_pair_shroud, Obsmetrologia.CausalSets.BlackHoles.Horizon.criticalCollapseCandidate_next_in_hole, Obsmetrologia.CausalSets.BlackHoles.Horizon.criticalCollapseCandidate_next_not_exterior, Obsmetrologia.CausalSets.BlackHoles.Horizon.criticalCollapseCandidate_holeRegionSwitch, Obsmetrologia.CausalSets.BlackHoles.Horizon.criticalCollapseCandidate_trappedSurfaceWitness, Obsmetrologia.CausalSets.BlackHoles.Horizon.MarkCriticalCollapseCandidate, Obsmetrologia.CausalSets.BlackHoles.Horizon.markCriticalCollapseCandidate_utterance_token, Obsmetrologia.CausalSets.BlackHoles.Horizon.markCriticalCollapseCandidate_utterance_claim, Obsmetrologia.CausalSets.BlackHoles.Horizon.markCriticalCollapseCandidate_tail_mem_prefixes, Obsmetrologia.CausalSets.BlackHoles.Horizon.markCriticalCollapseCandidate_next_mem_prefixes, Obsmetrologia.CausalSets.BlackHoles.Horizon.markCriticalCollapseCandidate_tail_reaches_next, Obsmetrologia.CausalSets.BlackHoles.Horizon.markCriticalCollapseCandidate_tail_causalLe_next, Obsmetrologia.CausalSets.BlackHoles.Horizon.markCriticalCollapseCandidate_tail_exterior, Obsmetrologia.CausalSets.BlackHoles.Horizon.markCriticalCollapseCandidate_next_shroud, Obsmetrologia.CausalSets.BlackHoles.Horizon.markCriticalCollapseCandidate_next_in_blackHoleRegion, Obsmetrologia.CausalSets.BlackHoles.Horizon.markCriticalCollapseCandidate_next_not_exterior}\leanok\uses{thm:kugelblitz-local-hole-entry, thm:finite-trapped-surface-witness}
The current finite critical-collapse carrier is not Choptuik scaling.  It is
the kugelblitz articulation step read at the formation threshold: the tail
endpoint is exterior, the next paired articulation lands on the shroud, hence
inside the no-readout region and not exterior, and the same data canonically
builds a finite trapped-surface witness.  The strengthened local readout also
records that the candidate tail and next point are the generated articulation
prefixes, that the tail reaches and causally precedes the next point, and that
the same laws specialize to an utterance assembled directly from token and
claim marks.
\end{theorem}

\begin{theorem}[External horizon-physics frontier map]\label{thm:external-horizon-physics-frontier-map}\notready\uses{thm:kugelblitz-local-hole-entry, thm:kugelblitz-growth-conjecture, thm:finite-so-readout-frame, thm:finite-apparent-event-plate, thm:finite-bekenstein-plate-counting, thm:finite-bekenstein-plate-capacity, thm:finite-horizon-membrane-dynamics, thm:readout-relative-horizon-frontier-laws, thm:finite-readout-patch-censorship-consolidation, thm:antihole-readout-complementarity, thm:finite-readout-no-comparison, thm:finite-wormhole-bridge, thm:finite-geon-carrier, thm:finite-ergosphere-extraction, thm:finite-unruh-readout-witness, thm:finite-unruh-thermal-readout-bridge, thm:finite-trapped-surface-witness, thm:finite-critical-collapse-carrier-laws}
The current corpus has constructed finite carriers for the black-hole-region formula,
apparent events, the counted shroud plate, finite Bekenstein-plate capacity as
positive Landauer support cost, readout-relative horizon packages, finite SO
readout frames,
shroud-indexed membrane dynamics, hole-region patches with no-exterior laws,
canonical finite censorship laws, finite antiholes with no listed-readout
ingress,
readout-complementarity witnesses with a pointwise no-comparison law, finite
trapped surface witnesses, and the finite wormhole bridge, ergosphere
extraction, Unruh readout witness, and finite thermal-readout bridge, plus the
kugelblitz/critical-collapse candidate
\lean{Obsmetrologia.CausalSets.BlackHoles.Horizon.CriticalCollapseCandidate}.
It has not yet proved the continuum trapped-surface singularity theorem,
Choptuik scaling, a continuum membrane paradigm, a continuum Unruh-effect
theorem,
continuum complementarity, smooth \(\mathrm{SO}(n)\) frame bundles, smooth
wormhole geometry, rotating continuum ergospheres, or continuum geons.  Those
remain red until the local carriers and laws are constructed without importing
the continuum vocabulary as explanation.
\end{theorem}

\begin{proposition}[Finite equation of state]\label{thm:finite-eos}\lean{Obsmetrologia.DenseMatter.Neutron.FiniteEOS}\leanok\uses{thm:finite-boltzmann-collision-lfp, thm:dft-minimizer-equilibrium}
A finite sampled equation of state records pressure, energy, and a finite sample
set.  It is a discrete admissibility projection, not neutron-star structure
equations.
\end{proposition}

\begin{theorem}[Admissibility has a speed bound]\label{thm:eos-speed-bound}\lean{Obsmetrologia.DenseMatter.Neutron.admissible_speed_bound}\leanok\uses{thm:finite-eos}
A stable causal finite equation of state satisfies the sampled pressure-speed
bound.  The stiff model is admissible and the superluminal model fails
causality, so the constraint has content.
\end{theorem}

\begin{proposition}[Nuclear-pasta excursion complex]\label{thm:pasta-cell-complex}\lean{Obsmetrologia.DenseMatter.Neutron.CellComplex}\leanok\uses{thm:cover-nerve, thm:persistent-h0-components-merge}
The pasta projection is a finite cell complex with vertices, edges, triangular
faces, and thresholded excursion subcomplexes.
\end{proposition}

\begin{theorem}[Excursion topology distinguishes filled and hollow triangles]\label{thm:pasta-euler-witnesses}\lean{Obsmetrologia.DenseMatter.Neutron.filledTriangle_euler}\leanok\uses{thm:pasta-cell-complex}
The filled triangle has Euler characteristic \(1\), while the companion
\lean{Obsmetrologia.DenseMatter.Neutron.hollowTriangle_euler} has Euler
characteristic \(0\).  This is finite excursion topology, not continuum nuclear
pasta morphology.
\end{theorem}

\begin{theorem}[Continuum dense-matter physics]\label{thm:continuum-dense-matter-frontier}\notready\uses{thm:eos-speed-bound, thm:pasta-euler-witnesses, thm:wilson-loop-invariant, thm:hole-future-closed}
The finite projections above do not yet prove TOV structure equations,
\(\chi\)EFT/nuclear many-body models, smooth event horizons, \(SU(3)\) gauge
theory, Wilson area laws, or confinement.  Those require the analytic
measure/Hilbert/PDE/Lie-bundle layers.
\end{theorem}
